\section{Scriptural Date and Location}
\label{sec:date-location}
{\small Each passage the Mass reads or sings from Scripture, in canonical order,
with what is dated, where it belongs, and the date the Scripture chronology corpus
gives for it. A passage whose parts are dated differently has a row for each part;
rows that share one answer name every element that uses it. The notes beneath give
each passage's setting, attribution and dating evidence.\par}
\medskip

\begin{chronologyorientationtable}
\chronologyorientationrow{row-18bc7e647866f638}{Written in Persia, as the article ``Esther'' in the Catholic Encyclopedia holds; the story is set at the Persian court in Susa (Esth 1:2), and Mardochai's prayer stands in the Greek chapters (13:8--18).}
\midrule
\chronologyorientationrow{row-d7ec5d94ad1d0e1d}{Place of writing unknown; the story is set in the land of Hus (Job 1:1), which Lam 4:21 joins to Edom.}
\midrule
\chronologyorientationrow{row-b12891d28424f424}{No psalm names an occasion or a place; the unsung title of Ps 89 names Moses, and Ps 118 names no author.}
\chronologyorientationrow{row-483ed5cea1be714a}{The Psalter as a collection; no individual psalm is dated, and no place of writing is named.}
\chronologyorientationrow{row-602ada67f8f3bb7d}{The psalm remembers Israel going out of Egypt and through the sea (Ps 113:1--3); the years are reckoned from the founding of the Temple (3 Kgs 6:1).}
\chronologyorientationrow{row-b5b7957d62c25282}{David, king in Jerusalem, to whom tradition attributes Ps 118; the reign dates David, not the psalm.}
\midrule
\chronologyorientationrow{row-88321469d09df574}{Jesus tells the parable to his disciples at Capharnaum in Galilee (Mt 17:23; 18:1); its king and servants are characters of a parable.}
\chronologyorientationrow{row-02ce0f03aad7eafe}{St Matthew writes in Palestine, by the tradition the Catholic Encyclopedia reports.}
\chronologyorientationrow{row-da30707e22c82488}{Writing place: Antioch in Syria, which the NABRE introduction offers as a plausible suggestion.}
\midrule
\chronologyorientationrow{row-d944994c29115eb5}{Paul writes in chains (Eph 6:20), at Rome or Caesarea, to Ephesus or to a group of churches.}
\chronologyorientationrow{row-d39c8c8473070f6e}{A later disciple developing Paul's thought, as one hypothesis the NABRE introduction offers; no place is given.}
\end{chronologyorientationtable}

\subsection{Esther 13:9--11, the Introit}
\label{sec:date-esther}

The Introit's prayer belongs to the chapters of Esther that the Vulgate prints
after the end of the Hebrew book (10:4--16:24), which St Jerome translated from the
Greek and set apart; the \work{Catholic Encyclopedia}'s ``Esther'' (vol.~5, 1909)
says that no trace of them is found in the Hebrew Bible and that they are most
probably translations from a Hebrew or Chaldaic original. The story is set at the
court of Assuerus, king of Persia, in Susa (1:2; 11:3); the prayer is said there,
in the days of Aman's decree. The date shown is that article's for the writing of
the book, placed at the end of the reign of Xerxes~I or under his son
Artaxerxes~I; the article calls it an early dating and allows a later one. The
chronology record applies the book's date to the Greek chapters as well, and it
dates no event of the story. The colophon at 11:1, which names the fourth year of a
Ptolemy and Cleopatra as the time this ``epistle of Phurim'' was brought, is the
book's own note on its Greek form and is not the date shown.

\subsection{Job 1--2, the Offertory}
\label{sec:date-job}

The Offertory narrates Job 1:1, 1:12, 1:13--19 and 2:6--7 in words of its own; the
chronology record answers every locus it cites with the same claim for the whole
book. The date shown is the modern critical judgment of the introduction to Job in
the \work{New American Bible, Revised Edition}, summarised there as a probable
period of writing; the record holds no traditional attribution and no date for
Job's life. St Gregory the Great reports that some held Moses to be the writer and
others one of the prophets, says that it is superfluous to ask who wrote it since
the Holy Spirit is its author, and then judges it the more probable that Job
himself, who bore the conflict, also told of his victory, speaking of himself as
of another, as Moses does (\work{Moralia}, Preface 1--3). The land of Hus, where the story is set,
is not the place of writing, which is unknown.

\subsection{Psalms 89, 113 and 118, the Gradual, Alleluia, Introit verse and Communion}
\label{sec:date-psalms}

None of the four psalm elements has an occasion in the record, and the
``Historical setting'' row says so for the Gradual, the Introit's verse and the
Communion. The composition date shown for all four is the bound that the
introduction to the Psalms in the \work{New American Bible, Revised Edition} sets for
the whole Psalter, within which, it says, no individual psalm can be dated
securely.

Ps~89 is titled \latin{Oratio Moysi, hominis Dei}, ``A prayer of Moses the man of
God'', a title the Gradual does not sing. The chronology record holds no
attribution for it, and no date is shown for Moses. The title names Moses. St
Augustine holds that the psalm was not ``entirely the composition of that Moses'',
since it lacks the expressions of Moses' songs (\work{Enarrationes} 89.2). St Robert
Bellarmine, in the general prefaces of his \work{Explanation of the Psalms}, admits
as acceptable, and in the preface that survives only in manuscript as not unlikely,
the opinion that those named in the titles, Moses among them, wrote their psalms;
in both he prefers as the more probable the opinion that David wrote all the
psalms, which he gives to St Augustine and St John Chrysostom. Neither Doctor
makes Moses the whole psalm's author; they differ only in that Bellarmine still
admits Moses' authorship as a tenable opinion. The exposition of the psalms printed under
St Athanasius's name says that Moses wrote the prayer and handed it on for those
who would come after (PG~27, col.~395). No witness read names an occasion or a
date.

Ps~113 names no author. The row for the Alleluia dates the event the psalm
remembers, the going out of Egypt, not the psalm: the preferred reckoning is the
Vulgate's statement that the Temple was begun in the four hundred and eightieth
year after the Exodus (3~Kgs 6:1), and the alternative, kept beside it, is the
Septuagint's figure, reported in the \work{Catholic Encyclopedia}'s ``Biblical
Chronology'' (vol.~3, 1908).

Ps~118 has no author in its title. The tradition attributes it to David: Haydock's
note at the first verse says that it is generally attributed to him and records
several proposed occasions without settling one. The reign shown is David's in the
``usual chronology'' of the \work{Catholic Encyclopedia}'s ``David'' (vol.~4, 1908),
marked disputed because the article reports a later dating beside it. It is the
reference point of the attribution and not a date of the psalm. The same rows serve
the Introit's verse, Ps~118:1, and the Communion's three half-verses.

\subsection{Matthew 18:23--35, the Gospel}
\label{sec:date-matthew}

The parable is told in the discourse of Mt~18, after the arrival at Capharnaum
(17:23) and before Jesus leaves Galilee (19:1). The record holds no date for that
discourse, and the ``Event'' row says so; the king, the servants and their debts
belong to a parable and have no date. The composition dates shown are the
\work{Catholic Encyclopedia}'s, from ``Gospel of Saint Matthew'' (vol.~10, 1911) and
``The New Testament'' (vol.~14, 1912), all marked disputed; the first of these
reports the tradition that Matthew wrote in Hebrew in Palestine, according to
Eusebius when he was leaving it.
The comparison row is the modern critical horizon of the introduction to Matthew in
the \work{New American Bible, Revised Edition}, which concerns the Greek Gospel; it is
shown beside the default answer and does not enter it. For the place where that
Gospel was composed, the introduction calls Antioch, the capital of the Roman
province of Syria, a plausible suggestion, not a settled fact.

\subsection{Ephesians 6:10--17, the Epistle}
\label{sec:date-ephesians}

Paul names himself (1:1) and writes as a prisoner (3:1; 4:1) and ``an ambassador in
a chain'' (6:20). The first range shown is that of the \work{Catholic
Encyclopedia}'s ``Epistle to the Ephesians'' (vol.~5, 1909), which sets the letter in
an imprisonment no earlier than the captivity at Caesarea and reports that the
greater number of interpreters take it as a circular letter; the year is from the
table in its ``St.\ Paul'' (vol.~11, 1911), which places the letter in the Roman
captivity with Philemon, Colossians and Philippians. Both are marked disputed. The
comparison row is the hypothesis that the introduction to Ephesians in the \work{New
American Bible, Revised Edition} offers as one possible explanation of the letter's
style and teaching, a later disciple developing Paul's thought; it is shown beside
the default answer and does not enter it.
